\section{Introduction}
This chapter discuss and dyamic programming algorithms for finding optimal policies.
Since these are dynamic prgramming algorithms, a perfect system model is required; this mean that
it is rarely used in practice, but these algorithms important foundations of the model-free algorithms.
Three algorithms will be presented;
\begin{itemize}
    \item Value Iteration: The algorithms suggested by the contraction mapping theorem \ref{thm:SolutionBellmanOptEq};
    \item Policy Iteration:
    \item Trucated Policy Iteration: The unified algorithm which includes value iteration and policy iteration as its special cases.
\end{itemize}

\section{Value Iteration}
    \subsection{Main Idea}
        The main idea of \textbf{value iteration} algorithms is the algorithms suggested in Theorem \ref{thm:SolutionBellmanOptEq}.
        In particular, the algorithms is
        \begin{equation}
            v_{k+1} = \max_{\pi \in \Pi}{(r_\pi + \gamma P_\pi v_k)},\quad k=0,1,2,\dots
        \end{equation}
        Breaking down, this iterative algorithm has two steps;
        \begin{enumerate}
            \item \emph{Policy Update}: The algorithm finds a policy $\pi_{k+1}$ that solves following optimization problem;
                \begin{equation}
                    \arg\max_\pi(r_\pi + \gamma P_\pi v_k)
                \end{equation}
            \item \emph{Value Update}: With the policy $\pi_{k+1}$, the algorithm calculates a new state value function $v_{k+1}$;
                \begin{equation}
                    v_{k+1} = r_{\pi_{k+1}} + \gamma P_{\pi_{k+1}} v_k
                \end{equation} and this $v_{k+1}$ is used in next iteration.
        \end{enumerate}
    \subsection{Pseudocode}
        \Algorithm{valueIter}
        {Value Iteration Algorithm}
        {\inputitem{$p(s',r|s,a)$}{the dynamics of the system} \\
         \inputitem{$v$}{initial guess}}
        {\inputitem{$v^*,\pi^*$}{optimal state value and optimal policy}}
        {
            $\pi \leftarrow$ greedy policy derived from $v$\;
            \Repeat{convergence}
            {
                \For{$s \in \mathcal{S}$} {
                    \For{$a \in \mathcal{A}(s)$} {
                        $q(s,a) \leftarrow \sum_r p(r|s,a)r + \gamma \sum_{s'}p(s'|s,a)v(s')$\;
                    }
                    $a^* \leftarrow \arg\max_a q(s,a)$\;
                    \tcp{policy update}
                    $\pi(s,a)=
                    \begin{cases}
                        1 &\text{if $a$ = $a^*$} \\
                        0 &\text{otherwise}
                    \end{cases}$\;
                    \tcp{value update}
                    $v(s) = \max_a{q(s,a)}$\;
                }
            }
            \Return{$v$, $\pi$}
        }
        Note that the value update of Algorithm \ref{alg:valueIter} is simply substituting the original state values with
        maximum action values since we are using greedy policy.
    \subsection{Convergence}
        The convergence can be proved with Theorem \ref{thm:ContractMap} and Theorem \ref{thm:ContPropBellmanOptEq}.
\section{Policy Iteration}
    \subsection{Main Idea}
        The \textbf{policy iteration} is based on the policy improvement theorem. 
        Like value iteration, the policy iteration consists of two steps:
        \begin{enumerate}
            \item \emph{Policy Evaluation}: For a given policy $\pi_k$, we calculate the corresponding state values by solving following Bellman equation:
                \begin{equation}
                    v_{\pi_k} = r_{\pi_k} + \gamma P_{\pi_k}v_{\pi_k}
                \end{equation}
            \item \emph{Policy Improvement}: Using the computed state value $v_k$, we obtain the improved policy as
                \begin{equation}
                    \pi_{k+1} = \arg\max_\pi(r_\pi + \gamma P_\pi v_{\pi_k})
                \end{equation} and here is where the policy improvement theorem is applied.
        \end{enumerate}
        The difference is that the value iteration is that instead of taking a state value and computing the policy from it,
        the policy iteration takes a policy and compute the value corresponding to the policy.

        Also, note that policy iteration has another iterative algorithm inside it: the policy evaluation uses iterative algorithm for solving the
        Bellman equation $v_{\pi_k} = r_{\pi_k} + \gamma P_{\pi_k}v_{\pi_k}$.

        